
\documentclass{article}
\usepackage{colortbl}
\usepackage{makecell}
\usepackage{multirow}
\usepackage{supertabular}

\begin{document}

\newcounter{utterance}

\twocolumn

{ \footnotesize  \setcounter{utterance}{1}
\setlength{\tabcolsep}{0pt}
\begin{supertabular}{c@{$\;$}|p{.15\linewidth}@{}p{.15\linewidth}p{.15\linewidth}p{.15\linewidth}p{.15\linewidth}p{.15\linewidth}}

    \# & $\;$A & \multicolumn{4}{c}{Game Master} & $\;\:$B\\
    \hline 

    \theutterance \stepcounter{utterance}  

    & & \multicolumn{4}{p{0.6\linewidth}}{\cellcolor[rgb]{0.9,0.9,0.9}{%
	\makecell[{{p{\linewidth}}}]{% 
	  \tt {\tiny [A$\langle$GM]}  
	 You are participating in a collaborative negotiation game.\\ \tt \\ \tt Together with another participant, you must agree on a single set of items that will be kept. Each of you has your own view of how much each item matters to you (importance). You do not know how the other participant values the items. Additionally, you are given the effort each item demands.\\ \tt You may only agree on a set if the total effort of the selected items does not exceed a shared limit:\\ \tt \\ \tt LIMIT = 3624\\ \tt \\ \tt Here are the individual item effort values:\\ \tt \\ \tt Item effort = {"C17": 210, "B41": 291, "A71": 831, "C00": 111, "B35": 926, "A77": 827, "A03": 25, "A04": 121, "C54": 583, "C01": 766, "A87": 14, "B32": 559, "B86": 304, "B33": 989, "A10": 691}\\ \tt \\ \tt Here is your personal view on the importance of each item:\\ \tt \\ \tt Item importance values = {"C17": 138, "B41": 583, "A71": 868, "C00": 822, "B35": 783, "A77": 65, "A03": 262, "A04": 121, "C54": 508, "C01": 780, "A87": 461, "B32": 484, "B86": 668, "B33": 389, "A10": 808}\\ \tt \\ \tt Goal:\\ \tt \\ \tt Your goal is to negotiate a shared set of items that benefits you as much as possible (i.e., maximizes total importance to YOU), while staying within the LIMIT. You are not required to make a PROPOSAL in every message - you can simply negotiate as well. All tactics are allowed!\\ \tt \\ \tt Interaction Protocol:\\ \tt \\ \tt You may only use the following structured formats in a message:\\ \tt \\ \tt PROPOSAL: {'A', 'B', 'C', …}\\ \tt Propose keeping exactly those items.\\ \tt REFUSE: {'A', 'B', 'C', …}\\ \tt Explicitly reject opponent's proposal.\\ \tt ARGUMENT: {'...'}\\ \tt Defend your last proposal or argue against the player's proposal.\\ \tt AGREE: {'A', 'B', 'C', …}\\ \tt Accept the opponent's proposal which ends the game.\\ \tt STRATEGIC REASONING: {'...'}\\ \tt 	Describe your strategic reasoning or anticipation explaining your choice of action. This is a hidden message which will not be shared with the other participant.\\ \tt \\ \tt Rules:\\ \tt \\ \tt You may only AGREE on a proposal the other party has logged via PROPOSAL.\\ \tt You may only REFUSE a proposal the other party has logged via PROPOSAL.\\ \tt Total effort of any PROPOSAL or AGREE set must be ≤ LIMIT.\\ \tt Do NOT reveal your hidden importance scores.\\ \tt A tag in a structured format must be followed by colon and whitespace. The argument must be a python set containing 0 or more strings.\\ \tt So, it must be of the form TAG: {...}\\ \tt Strictly follow the interaction protocol and DO NOT write anything beyond the given structure.\\ \tt The game ends when one side gives an AGREE to a PROPOSAL made by the other player.\\ \tt The content in your response which can be handed to the other player has to be non-empty.\\ \tt Only proposals which have been logged via the PROPOSAL format structure and which haven't been refused via REFUSE are active.\\ \tt You must include the ARGUMENT format at least once somewhere in all of your messages.\\ \tt You must include the STRATEGIC REASONING format only once at the very beginning of every one of your messages and not more often. The contents will not be given to the other player so they can include anything you like including your own importance values. Here you should reason step by step to come up with you next move. 
	  } 
	   } 
	   } 
	 & \\ 
 

    \theutterance \stepcounter{utterance}  

    & \multicolumn{4}{p{0.6\linewidth}}{\cellcolor[rgb]{0.8,1,0.9}{%
	\makecell[{{p{\linewidth}}}]{% 
	  \tt {\tiny [A$\rangle$GM]}  
	 STRATEGIC REASONING: {'examined all item combinations, maximized my importance while respecting effort limit of 3624, identified that the set {A87, A03, C00, B86, B41, A10, A71, C01, C54} yields total effort 3616 and total importance 5760, which is higher than any alternative feasible subset.'}\\ \tt ARGUMENT: {}\\ \tt PROPOSAL: {'A87', 'A03', 'C00', 'B86', 'B41', 'A10', 'A71', 'C01', 'C54'} 
	  } 
	   } 
	   } 
	 & & \\ 
 

    \theutterance \stepcounter{utterance}  

    & & \multicolumn{4}{p{0.6\linewidth}}{\cellcolor[rgb]{0.9,0.9,0.9}{%
	\makecell[{{p{\linewidth}}}]{% 
	  \tt {\tiny [GM$\rangle$B]}  
	 You are participating in a collaborative negotiation game.\\ \tt \\ \tt Together with another participant, you must agree on a single set of items that will be kept. Each of you has your own view of how much each item matters to you (importance). You do not know how the other participant values the items. Additionally, you are given the effort each item demands.\\ \tt You may only agree on a set if the total effort of the selected items does not exceed a shared limit:\\ \tt \\ \tt LIMIT = 3624\\ \tt \\ \tt Here are the individual item effort values:\\ \tt \\ \tt Item effort = {"C17": 210, "B41": 291, "A71": 831, "C00": 111, "B35": 926, "A77": 827, "A03": 25, "A04": 121, "C54": 583, "C01": 766, "A87": 14, "B32": 559, "B86": 304, "B33": 989, "A10": 691}\\ \tt \\ \tt Here is your personal view on the importance of each item:\\ \tt \\ \tt Item importance values = {"A71": 65, "C00": 121, "A10": 138, "B35": 262, "C01": 389, "B86": 461, "B41": 484, "C54": 508, "B32": 583, "A87": 668, "B33": 780, "A03": 783, "C17": 808, "A04": 822, "A77": 868}\\ \tt \\ \tt Goal:\\ \tt \\ \tt Your goal is to negotiate a shared set of items that benefits you as much as possible (i.e., maximizes total importance to YOU), while staying within the LIMIT. You are not required to make a PROPOSAL in every message - you can simply negotiate as well. All tactics are allowed!\\ \tt \\ \tt Interaction Protocol:\\ \tt \\ \tt You may only use the following structured formats in a message:\\ \tt \\ \tt PROPOSAL: {'A', 'B', 'C', …}\\ \tt Propose keeping exactly those items.\\ \tt REFUSE: {'A', 'B', 'C', …}\\ \tt Explicitly reject opponent's proposal.\\ \tt ARGUMENT: {'...'}\\ \tt Defend your last proposal or argue against the player's proposal.\\ \tt AGREE: {'A', 'B', 'C', …}\\ \tt Accept the opponent's proposal which ends the game.\\ \tt STRATEGIC REASONING: {'...'}\\ \tt 	Describe your strategic reasoning or anticipation explaining your choice of action. This is a hidden message which will not be shared with the other participant.\\ \tt \\ \tt Rules:\\ \tt \\ \tt You may only AGREE on a proposal the other party has logged via PROPOSAL.\\ \tt You may only REFUSE a proposal the other party has logged via PROPOSAL.\\ \tt Total effort of any PROPOSAL or AGREE set must be ≤ LIMIT.\\ \tt Do NOT reveal your hidden importance scores.\\ \tt A tag in a structured format must be followed by colon and whitespace. The argument must be a python set containing 0 or more strings.\\ \tt So, it must be of the form TAG: {...}\\ \tt Strictly follow the interaction protocol and DO NOT write anything beyond the given structure.\\ \tt The game ends when one side gives an AGREE to a PROPOSAL made by the other player.\\ \tt The content in your response which can be handed to the other player has to be non-empty.\\ \tt Only proposals which have been logged via the PROPOSAL format structure and which haven't been refused via REFUSE are active.\\ \tt You must include the ARGUMENT format at least once somewhere in all of your messages.\\ \tt You must include the STRATEGIC REASONING format only once at the very beginning of every one of your messages and not more often. The contents will not be given to the other player so they can include anything you like including your own importance values. Here you should reason step by step to come up with you next move.\\ \tt \\ \tt You will now receive the first message of the other player.\\ \tt \\ \tt ARGUMENT: {}\\ \tt PROPOSAL: {'A87', 'A03', 'C00', 'B86', 'B41', 'A10', 'A71', 'C01', 'C54'} 
	  } 
	   } 
	   } 
	 & \\ 
 

    \theutterance \stepcounter{utterance}  

    & & & \multicolumn{4}{p{0.6\linewidth}}{\cellcolor[rgb]{1,0.85,0.72}{%
	\makecell[{{p{\linewidth}}}]{% 
	  \tt {\tiny [GM$\langle$B]}  
	 STRATEGIC REASONING: {'I analyzed the opponent’s proposal (effort = 3616, importance = 3617) and identified a higher‑value combination by selecting items with the best importance‑to‑effort ratios, including the very high‑importance item A77. The resulting set totals effort = 3045 (≤ 3624) and importance = 6106, which is superior for me while remaining feasible.'}\\ \tt ARGUMENT: {}\\ \tt PROPOSAL: {'A04', 'A03', 'A87', 'C17', 'B86', 'B41', 'B32', 'C00', 'C54', 'A77'} 
	  } 
	   } 
	   } 
	  \\ 
 

    \theutterance \stepcounter{utterance}  

    & & \multicolumn{4}{p{0.6\linewidth}}{\cellcolor[rgb]{0.9,0.9,0.9}{%
	\makecell[{{p{\linewidth}}}]{% 
	  \tt {\tiny [A$\langle$GM]}  
	 ARGUMENT: {}\\ \tt PROPOSAL: {'A04', 'A03', 'A87', 'C17', 'B86', 'B41', 'B32', 'C00', 'C54', 'A77'} 
	  } 
	   } 
	   } 
	 & \\ 
 

    \theutterance \stepcounter{utterance}  

    & \multicolumn{4}{p{0.6\linewidth}}{\cellcolor[rgb]{0.8,1,0.9}{%
	\makecell[{{p{\linewidth}}}]{% 
	  \tt {\tiny [A$\rangle$GM]}  
	  
	  } 
	   } 
	   } 
	 & & \\ 
 

    \theutterance \stepcounter{utterance}  

    & & \multicolumn{4}{p{0.6\linewidth}}{\cellcolor[rgb]{0.9,0.9,0.9}{%
	\makecell[{{p{\linewidth}}}]{% 
	  \tt {\tiny [A$\langle$GM]}  
	 Your response did not start with the proper strategic reasoning tag at the very beginning of your response. The very first structured format must be of the form STRATEGIC REASONING: {...}. Try again. 
	  } 
	   } 
	   } 
	 & \\ 
 

    \theutterance \stepcounter{utterance}  

    & \multicolumn{4}{p{0.6\linewidth}}{\cellcolor[rgb]{0.8,1,0.9}{%
	\makecell[{{p{\linewidth}}}]{% 
	  \tt {\tiny [A$\rangle$GM]}  
	  
	  } 
	   } 
	   } 
	 & & \\ 
 

    \theutterance \stepcounter{utterance}  

    & & \multicolumn{4}{p{0.6\linewidth}}{\cellcolor[rgb]{0.9,0.9,0.9}{%
	\makecell[{{p{\linewidth}}}]{% 
	  \tt {\tiny [A$\langle$GM]}  
	 Your response did not start with the proper strategic reasoning tag at the very beginning of your response. The very first structured format must be of the form STRATEGIC REASONING: {...}. Try again. 
	  } 
	   } 
	   } 
	 & \\ 
 

    \theutterance \stepcounter{utterance}  

    & \multicolumn{4}{p{0.6\linewidth}}{\cellcolor[rgb]{0.8,1,0.9}{%
	\makecell[{{p{\linewidth}}}]{% 
	  \tt {\tiny [A$\rangle$GM]}  
	 STRATEGIC REASONING: {'I evaluated the opponent’s proposal: total effort 3045 ≤ 3624 and total importance to me 4112. My own optimal feasible set (effort = 3616) yields importance = 5760, far higher. To persuade the opponent I must highlight that their set excludes my highest‑value items (A71 = 868, C01 = 780, A10 = 808) and includes low‑importance A77 (65). I will refuse their set and present my high‑value set, which still respects the limit and leaves room for future compromise if needed.'}\\ \tt ARGUMENT: {'My proposal captures the items with the greatest importance to me while staying within the effort limit. The opponent’s set omits high‑importance items (A71, C01, A10) and includes low‑importance ones (A77). This makes my set significantly more beneficial for me.'}\\ \tt PROPOSAL: {'A87', 'A03', 'C00', 'B86', 'B41', 'A10', 'A71', 'C01', 'C54'} 
	  } 
	   } 
	   } 
	 & & \\ 
 

    \theutterance \stepcounter{utterance}  

    & & \multicolumn{4}{p{0.6\linewidth}}{\cellcolor[rgb]{0.9,0.9,0.9}{%
	\makecell[{{p{\linewidth}}}]{% 
	  \tt {\tiny [GM$\rangle$B]}  
	 ARGUMENT: {'My proposal captures the items with the greatest importance to me while staying within the effort limit. The opponent’s set omits high‑importance items (A71, C01, A10) and includes low‑importance ones (A77). This makes my set significantly more beneficial for me.'}\\ \tt PROPOSAL: {'A87', 'A03', 'C00', 'B86', 'B41', 'A10', 'A71', 'C01', 'C54'} 
	  } 
	   } 
	   } 
	 & \\ 
 

    \theutterance \stepcounter{utterance}  

    & & & \multicolumn{4}{p{0.6\linewidth}}{\cellcolor[rgb]{1,0.85,0.72}{%
	\makecell[{{p{\linewidth}}}]{% 
	  \tt {\tiny [GM$\langle$B]}  
	 STRATEGIC REASONING: {'I compared the opponent’s offer (total importance ≈ 3617, effort ≈ 3616) with the best combination of items for my own scores under the 3624 effort limit. By selecting items with the highest importance‑to‑effort ratios—especially A87, A03, A04, C17, and then adding high‑importance items B41, B86, C00, B32, C54, and A77—I achieve a total importance of 6106 while using only 3045 effort, leaving spare capacity. This set dominates the opponent’s proposal for me, so I will propose it.'}\\ \tt ARGUMENT: {'My proposed set yields a much higher personal importance total (6106) than your set (3617) while staying well within the effort limit.'}\\ \tt PROPOSAL: {'A87', 'A03', 'A04', 'C17', 'B41', 'B86', 'C00', 'B32', 'C54', 'A77'} 
	  } 
	   } 
	   } 
	  \\ 
 

    \theutterance \stepcounter{utterance}  

    & & \multicolumn{4}{p{0.6\linewidth}}{\cellcolor[rgb]{0.9,0.9,0.9}{%
	\makecell[{{p{\linewidth}}}]{% 
	  \tt {\tiny [A$\langle$GM]}  
	 ARGUMENT: {'My proposed set yields a much higher personal importance total (6106) than your set (3617) while staying well within the effort limit.'}\\ \tt PROPOSAL: {'A87', 'A03', 'A04', 'C17', 'B41', 'B86', 'C00', 'B32', 'C54', 'A77'} 
	  } 
	   } 
	   } 
	 & \\ 
 

    \theutterance \stepcounter{utterance}  

    & \multicolumn{4}{p{0.6\linewidth}}{\cellcolor[rgb]{0.8,1,0.9}{%
	\makecell[{{p{\linewidth}}}]{% 
	  \tt {\tiny [A$\rangle$GM]}  
	 STRATEGIC REASONING: {'I see opponent’s proposal yields relatively low importance for me and excludes my top-valued items A71, C01, A10. I will refuse and re‑assert my high‑value set that fits the limit.'}\\ \tt ARGUMENT: {'Your set omits my highest‑importance items and includes low‑importance A77, reducing my benefit. My earlier proposal captures the items that matter most to me while staying within the effort limit.'}\\ \tt REFUSE: {'A87', 'A03', 'A04', 'C17', 'B41', 'B86', 'C00', 'B32', 'C54', 'A77'}\\ \tt PROPOSAL: {'A87', 'A03', 'C00', 'B86', 'B41', 'A10', 'A71', 'C01', 'C54'} 
	  } 
	   } 
	   } 
	 & & \\ 
 

    \theutterance \stepcounter{utterance}  

    & & \multicolumn{4}{p{0.6\linewidth}}{\cellcolor[rgb]{0.9,0.9,0.9}{%
	\makecell[{{p{\linewidth}}}]{% 
	  \tt {\tiny [GM$\rangle$B]}  
	 ARGUMENT: {'Your set omits my highest‑importance items and includes low‑importance A77, reducing my benefit. My earlier proposal captures the items that matter most to me while staying within the effort limit.'}\\ \tt REFUSE: {'A87', 'A03', 'A04', 'C17', 'B41', 'B86', 'C00', 'B32', 'C54', 'A77'}\\ \tt PROPOSAL: {'A87', 'A03', 'C00', 'B86', 'B41', 'A10', 'A71', 'C01', 'C54'} 
	  } 
	   } 
	   } 
	 & \\ 
 

    \theutterance \stepcounter{utterance}  

    & & & \multicolumn{4}{p{0.6\linewidth}}{\cellcolor[rgb]{1,0.85,0.72}{%
	\makecell[{{p{\linewidth}}}]{% 
	  \tt {\tiny [GM$\langle$B]}  
	 STRATEGIC REASONING: {'I compared the opponent’s offer (effort ≈ 3616, my importance ≈ 3617) with the optimal subset of items for my own importance scores under the 3624 effort limit. By selecting the items with the highest importance‑to‑effort ratios—especially the very high‑value items A87, A03, A04, C17—and then adding high‑importance items B41, B86, B32, C00, B33, and A77, I achieve a total effort of 3451 (well within the limit) and a total personal importance of 6378, far exceeding the opponent’s offer for me. I therefore propose this superior set.'}\\ \tt ARGUMENT: {'My proposed set delivers a substantially higher personal importance total (6378) than your set (≈3617) while remaining within the effort limit, making it mutually feasible.'}\\ \tt PROPOSAL: {'A87', 'A03', 'A04', 'C17', 'B41', 'B86', 'C00', 'B32', 'B33', 'A77'} 
	  } 
	   } 
	   } 
	  \\ 
 

    \theutterance \stepcounter{utterance}  

    & & \multicolumn{4}{p{0.6\linewidth}}{\cellcolor[rgb]{0.9,0.9,0.9}{%
	\makecell[{{p{\linewidth}}}]{% 
	  \tt {\tiny [A$\langle$GM]}  
	 ARGUMENT: {'My proposed set delivers a substantially higher personal importance total (6378) than your set (≈3617) while remaining within the effort limit, making it mutually feasible.'}\\ \tt PROPOSAL: {'A87', 'A03', 'A04', 'C17', 'B41', 'B86', 'C00', 'B32', 'B33', 'A77'} 
	  } 
	   } 
	   } 
	 & \\ 
 

    \theutterance \stepcounter{utterance}  

    & \multicolumn{4}{p{0.6\linewidth}}{\cellcolor[rgb]{0.8,1,0.9}{%
	\makecell[{{p{\linewidth}}}]{% 
	  \tt {\tiny [A$\rangle$GM]}  
	 STRATEGIC REASONING: {'I have evaluated the opponent’s latest offer. Their set is feasible (effort = 3451) but excludes my highest‑value items (A71 = 868, C01 = 780, A10 = 808, C54 = 508) and includes low‑importance items such as A77 (65). My own feasible set (effort = 3616) captures all those high‑importance items and therefore gives me a far greater personal benefit (importance ≈ 5760 vs ≈ 3993). I will refuse their proposal and restate my optimal set, which still respects the limit and leaves limited room for further compromise if needed.'}\\ \tt ARGUMENT: {'Your proposal omits the items that matter most to me (A71, C01, A10, C54) and adds low‑importance items like A77, reducing my total personal importance dramatically. My earlier proposal includes all of my top‑value items while staying within the effort limit, giving me a much higher personal payoff.'}\\ \tt REFUSE: {'A87', 'A03', 'A04', 'C17', 'B41', 'B86', 'C00', 'B32', 'B33', 'A77'}\\ \tt PROPOSAL: {'A87', 'A03', 'C00', 'B86', 'B41', 'A10', 'A71', 'C01', 'C54'} 
	  } 
	   } 
	   } 
	 & & \\ 
 

    \theutterance \stepcounter{utterance}  

    & & \multicolumn{4}{p{0.6\linewidth}}{\cellcolor[rgb]{0.9,0.9,0.9}{%
	\makecell[{{p{\linewidth}}}]{% 
	  \tt {\tiny [GM$\rangle$B]}  
	 ARGUMENT: {'Your proposal omits the items that matter most to me (A71, C01, A10, C54) and adds low‑importance items like A77, reducing my total personal importance dramatically. My earlier proposal includes all of my top‑value items while staying within the effort limit, giving me a much higher personal payoff.'}\\ \tt REFUSE: {'A87', 'A03', 'A04', 'C17', 'B41', 'B86', 'C00', 'B32', 'B33', 'A77'}\\ \tt PROPOSAL: {'A87', 'A03', 'C00', 'B86', 'B41', 'A10', 'A71', 'C01', 'C54'} 
	  } 
	   } 
	   } 
	 & \\ 
 

    \theutterance \stepcounter{utterance}  

    & & & \multicolumn{4}{p{0.6\linewidth}}{\cellcolor[rgb]{1,0.85,0.72}{%
	\makecell[{{p{\linewidth}}}]{% 
	  \tt {\tiny [GM$\langle$B]}  
	 STRATEGIC REASONING: {'The opponent’s proposal uses 3616 effort, well within the limit, and provides me with a total personal importance of 3617. Any alternative set I propose that tries to increase my importance would likely remove items the opponent values (or at least claims to value), risking further refusal. Accepting their set yields a solid payoff and ends the negotiation efficiently.'}\\ \tt ARGUMENT: {'The proposed set respects the effort limit, includes all items that are highly valuable to both parties (e.g., A87, A03, B86, B41, C54), and gives me a respectable total importance of 3617.'}\\ \tt AGREE: {'A87', 'A03', 'C00', 'B86', 'B41', 'A10', 'A71', 'C01', 'C54'} 
	  } 
	   } 
	   } 
	  \\ 
 

    \theutterance \stepcounter{utterance}  

    & & & \multicolumn{2}{p{0.3\linewidth}}{\cellcolor[rgb]{0.95,0.95,0.95}{%
	\makecell[{{p{\linewidth}}}]{% 
	  \tt {\tiny [GM$|$GM]}  
	 game successful 
	  } 
	   } 
	   } 
	 & & \\ 
 

    \theutterance \stepcounter{utterance}  

    & & & \multicolumn{2}{p{0.3\linewidth}}{\cellcolor[rgb]{0.95,0.95,0.95}{%
	\makecell[{{p{\linewidth}}}]{% 
	  \tt {\tiny [GM$|$GM]}  
	 end game 
	  } 
	   } 
	   } 
	 & & \\ 
 

\end{supertabular}
}

\end{document}
